\section{Complete services in the original investment economy}\label{sec:study65}
The numerical question is whether the organization of a learned query changes the work of delivering an accurate policy. We use the original scalar investment primitives, discount factor $15/16$, and continuous uniform innovations throughout. The returned object is the complete callable controller: at each subsequently acquired state it obtains a fresh certificate before implementing its action. A finite workload measures its execution. The uniform original-law theorem, rather than an extrapolation from that workload, supplies its accuracy guarantee.

\subsection{Design, controls and measured objects}
Two protocols were frozen before their respective executions. The first compared adaptive parabolic search, bisection, unconditional ReLU insertion and unconditional quadratic insertion. The second retained those controls and added certificate-gated routing for each fitted predictor. Its insertion and routing services independently reconstructed exactly the same fitted model for a matched seed, task and repetition. The initial and innovation workloads were also identical within each task. This design separates query organization from a change in model parameters or evaluated states.

The task catalogue is $(d,T,N)=(2,2,128),(8,2,128),(16,2,128),(8,3,32)$, with original-optimum targets $1/32$ and $1/128$. Here $N$ is the number of workload trajectories. Each learned method uses both specified fitting seeds, and every service has two isolated-process repetitions. The first protocol therefore contains 96 services and the second 160. The fitting seeds are 6301 and 6302 in the first catalogue and 6401 and 6402 in the second. No successful seed, task, target or repetition is selected for publication.

Each fitted service pays for its own 96 training states, nested Bellman-generated labels, feature construction and fit. The ReLU predictor has one hidden layer of width 24 and uses a capped L-BFGS fit. The quadratic predictor uses the same features, their second-order products and regularization. The permitted training iterations and function evaluations are fixed in the deposited protocols. At two dates the training account records 6,912 Bellman action queries; at three dates it records 155,904. These costs are not assigned to a previously solved conventional state lattice and are not omitted when a boundary certificate avoids a later prediction.

The adaptive comparator chooses an unresolved interval using its parabolic minorant and a safeguarded proposed split. Bisection supplies an independent refinement rule. Both choose their own action partitions and stopping stages; all methods use the same analytic terminal kernel, original-law child quadrature and precision guard. A fitted model cannot relax these requirements. These are structure-exploiting conventional controls for the scalar action problem, not a sparse-grid package relabeled as a neural method.

Every process is pinned to one CPU, with one numerical-library thread. The parent clock starts before launching the child and stops after process exit and join. It includes imports, training labels, fitting, all deployment queries, exact workload-cost reconstruction, serialized arrays, durable summaries and the final clock record. Processor frequency is uncontrolled. We report individual repetitions and machine-independent query counts alongside clocks; small clock differences are not treated as statistically identified runtime effects.

\subsection{Complete time and query attribution}
Table~\ref{tab:clocks65} reports the second catalogue's full cold-start process times. A conventional method has the smallest median in every task--target cell: bisection in three cells and adaptive parabolic search in five. None of the learned services has a smaller median complete time. The repeated reconstruction of a fitted model is part of this cold-start estimand. A different reuse contract would require a different work account, not deletion of the recorded fitting charge.

\begin{table}[htbp]
\centering\caption{Complete cold-start service times}\label{tab:clocks65}
{\small\setlength{\tabcolsep}{4pt}
\begin{tabular}{lrrrrrrrr}
\toprule
$d$ & $T$ & $\epsilon$ & A & B & RI & RG & QI & QG \\
\midrule
2 & 2 & 1/32 & 0.256 & 0.252 & 0.650 & 0.657 & 0.614 & 0.614 \\
2 & 2 & 1/128 & 0.270 & 0.282 & 0.655 & 0.656 & 0.624 & 0.628 \\
8 & 2 & 1/32 & 0.380 & 0.360 & 0.770 & 0.771 & 0.741 & 0.739 \\
8 & 2 & 1/128 & 0.416 & 0.441 & 0.806 & 0.806 & 0.754 & 0.773 \\
16 & 2 & 1/32 & 0.536 & 0.496 & 0.935 & 0.911 & 0.891 & 0.856 \\
16 & 2 & 1/128 & 0.552 & 0.622 & 0.978 & 0.987 & 0.913 & 0.935 \\
8 & 3 & 1/32 & 0.984 & 1.031 & 1.962 & 1.813 & 1.840 & 1.802 \\
8 & 3 & 1/128 & 1.450 & 1.952 & 2.435 & 2.481 & 2.323 & 2.285 \\
\bottomrule
\end{tabular}}
\par\smallskip{\footnotesize A: adaptive parabolic search; B: bisection; RI and RG: ReLU insertion and gating; QI and QG: quadratic insertion and gating. Seconds are parent-observed complete-process times, including imports, training, queries, exact workload costing, durable output and process join. Each conventional median uses two repetitions; each learned median uses two seeds and two repetitions. Single-host descriptive clocks are not significance tests. All individual observations are retained in the supplement.}
\end{table}

\begin{table}[htbp]
\centering\caption{Deployment Bellman query counts on identical workloads}\label{tab:queries65}
{\small\setlength{\tabcolsep}{4pt}
\begin{tabular}{rrrrrrrrr}
\toprule
$d$ & $T$ & $\epsilon$ & A & B & RI & RG & QI & QG \\
\midrule
2 & 2 & 1/32 & 1,298 & 1,310 & 1,427 & 1,298 & 1,421 & 1,295 \\
2 & 2 & 1/128 & 2,508 & 2,658 & 2,618 & 2,433 & 2,623 & 2,468 \\
8 & 2 & 1/32 & 1,280 & 1,325 & 1,472 & 1,278.5 & 1,461.5 & 1,278.5 \\
8 & 2 & 1/128 & 2,373 & 2,648 & 2,613 & 2,375.5 & 2,588 & 2,365.5 \\
16 & 2 & 1/32 & 1,286 & 1,328 & 1,455.5 & 1,281.5 & 1,457 & 1,283 \\
16 & 2 & 1/128 & 2,318 & 2,633 & 2,578 & 2,358 & 2,603 & 2,360.5 \\
8 & 3 & 1/32 & 7,042 & 8,112 & 9,014 & 6,947.5 & 8,604 & 7,038.5 \\
8 & 3 & 1/128 & 31,015 & 39,768 & 37,519 & 31,074 & 35,295.5 & 30,898.5 \\
\bottomrule
\end{tabular}}
\par\smallskip{\footnotesize Counts exclude training and are not a complete work measure. They include the recursive descendants and the common terminal kernel. Within a task and tolerance the workload is fixed across methods. Half-integer medians arise from two different seed transcripts. Insertion and gating share identical fitted parameters for each matched seed and repetition.}
\end{table}


The deployment counts in Table~\ref{tab:queries65} answer a narrower question. Holding each fitted model fixed, routing avoids the unconditional interior insertion and instead uses a prediction only when a certificate requires another query. There are sixteen distinct task--target--seed comparisons for each predictor. Routing reduces deployment Bellman queries in all sixteen ReLU comparisons and all sixteen quadratic comparisons. The repeated traces are identical, so repetitions do not turn these 32 matched comparisons into 64 independent numerical outcomes.

Against adaptive parabolic search the result is mixed. The ReLU route uses fewer queries in nine comparisons, the same number in three, and more in four. The quadratic route uses fewer in eleven, the same in two, and more in three. In the two-state, two-date, $1/128$ cell, the median ReLU deployment count falls from 2,508 for adaptive search to 2,433 for routing. In the sixteen-state counterpart it rises from 2,318 to 2,358. The three-date, tighter-target ReLU median also exceeds the adaptive count, despite improving substantially on unconditional insertion. The gate does not claim a pathwise query advantage over a strong conventional split rule.

\begin{table}[htbp]
\centering\caption{Matched query differences: gated minus comparator}\label{tab:matched65}
{\small\setlength{\tabcolsep}{4pt}
\begin{tabular}{llrrr}
\toprule
Predictor & Comparator & Lower & Equal & Higher \\
\midrule
Relu & Insertion & 16 & 0 & 0 \\
Relu & Adaptive & 9 & 3 & 4 \\
Quadratic & Insertion & 16 & 0 & 0 \\
Quadratic & Adaptive & 11 & 2 & 3 \\
\bottomrule
\end{tabular}}
\par\smallskip{\footnotesize One row comparison per task, target and fitting seed: sixteen for each predictor and comparator. Repetitions have identical models and numerical transcripts and do not double this count. These are query-count classifications, not signs of expected policy-cost differences.}
\end{table}


These observations identify an incremental contribution of query organization, not a full economic dominance result. A prediction may reduce the upper endpoint, improve the lower certificate through localization, or do neither. Proposition~\ref{prop:localization65} describes the sufficient mathematical conditions; the matched catalogue measures their algorithmic consequence without giving the predictor credit for the certificate's common components. Fitting and model evaluation can exceed the value of the saved queries. Here the complete clocks show that they do.

\subsection{Accuracy, dimension and finite-catalogue reliability}
All 256 declared services returned their required local certificates and uniform policy allowance. The second catalogue contains twenty services at each task--target cell. Its strongest state-dimensional execution is $d=16$, $T=2$, with target $1/128$ and no stored state-lattice nodes. This is different from the earlier eight-dimensional tensor construction, which attained only the coarse target $1/2$. The economic primitives and original optimum are unchanged; the computation now queries descendants instead of filling a state grid.

\begin{table}[htbp]
\centering\caption{Accuracy, memory and finite-catalogue return}\label{tab:reliability65}
{\small\setlength{\tabcolsep}{4pt}
\begin{tabular}{rrrrrrrr}
\toprule
$d$ & $T$ & Target & Returns & Gap bound & Batch & RSS (KiB) & Max (s) \\
\midrule
2 & 2 & 1/32 & 20 & 0.01562501 & 256 & 94520 & 0.945 \\
2 & 2 & 1/128 & 20 & 0.00390626 & 512 & 94496 & 0.670 \\
8 & 2 & 1/32 & 20 & 0.01562501 & 256 & 94748 & 0.792 \\
8 & 2 & 1/128 & 20 & 0.00390626 & 512 & 95160 & 0.824 \\
16 & 2 & 1/32 & 20 & 0.01562501 & 256 & 96040 & 0.953 \\
16 & 2 & 1/128 & 20 & 0.00390626 & 512 & 96212 & 1.006 \\
8 & 3 & 1/32 & 20 & 0.01562501 & 512 & 95704 & 1.964 \\
8 & 3 & 1/128 & 20 & 0.00390626 & 2048 & 98788 & 2.583 \\
\bottomrule
\end{tabular}}
\par\smallskip{\footnotesize Every cell has twenty isolated-process returns: two conventional methods with two repetitions and four learned methods with two seeds and two repetitions. Gap bounds are rounded upward for display; exact binary64 values are retained in the audit and checked against the rational target. RSS is the largest recorded process peak, not only model storage. No service built a stored state lattice or invoked exact-rational recovery.}
\end{table}


The uniform bound is approximately half the stated target plus the explicit acquisition allowance, because the fixed design reserves half of the target when selecting the per-date local request. Every exact recorded bound, not its printed decimal, is checked against its rational target. The proof covers true initial states outside the workload. The executed workload alone would not establish that assertion.

Removing a state lattice does not make the problem dimensionless. Each original cost and transition reads all coordinates, the dense quadratic model has a growing feature array, and the batched continuation recursion retains innovation descendants. Across both catalogues, the largest live batch is 2,048 states and the largest process peak is 98,860 KiB. The three-date cells reveal the continuation-tree cost; the tight adaptive deployment count rises to 31,015 queries for only 32 trajectories. No tensor-free two-control or high-dimensional shock execution is inferred from these scalar-control, scalar-innovation results. The earlier complete two-control policy study remains in the preserved article under its own implementation.

Reliability is defined for the frozen finite catalogue. All services returned; the worst recorded complete process time is 2.5831 seconds when rounded upward to four decimals. There were 104 fitting warnings across the two catalogues, all retained. A fitting warning does not forge a returned Bellman certificate, but it also cannot be discarded from a training-reliability description. No floating query invoked the separately counted exact-rational recovery on these workloads. That observation does not eliminate the recovery procedure from the uniform implementation contract or assign zero cost to its possible use elsewhere.

\subsection{Policy cost and the purchase decision}
The workload starts from fixed twenty-bit dyadic states and uses fixed dyadic shocks, with eight named boundary initial states. Every true state transition and path cost is reconstructed with rational arithmetic. There are 26,624 path rows and 55,296 implemented decisions across the two catalogues. These are finite diagnostic workloads, not independent draws furnishing new confidence intervals under the continuous law. The rational workload mean is retained for exact reproduction and is not substituted for the expectation in Section~\ref{sec:purchase64}.

The preceding continuous-law policy comparisons consequently remain important and separate. All twenty pure fitted ReLU policies in the tensor study have strictly higher expected cost than the common policy; the twenty common-augmented comparisons are unresolved. The common policy has the smallest complete release work in all twenty attainable task--target cells. The query study neither overwrites those observations nor asserts a new expected-cost reversal on the basis of finite dyadic averages. A callable query controller, a pure fitted actor, and a common-augmented nodal policy are distinct objects.

The purchase criterion has two uses. A buyer imposing only a uniform accuracy constraint can choose the least expensive eligible observed service for the declared workload and hardware. This favors a conventional method in the current cold-start catalogue. A buyer comparing monetary expected gains must additionally supply a valid expected-gain interval and the numeraire conversion, computation price and fixed fee in Section~\ref{sec:purchase64}. Neither the current clocks nor two regret upper bounds provide that directional gain interval. The criterion is a specified decision problem, not an empirical calibration of an installation charge.

\subsection{Reproduction and preserved evidence}
The revision checks every service request, source freeze, model identity, trace, exact cost and complete-process receipt. An independently written rational implementation reconstructs the economic trajectories. Re-querying each distinct recorded observation reproduces the original numerical endpoints, actions, probe counts and operation ledger; byte-identical repetitions are separately checked by identity rather than counted as fresh simulations. This is more than checking a table, but less than a formal proof-checker verification of arbitrary interval code.

The present integration creates no new training run, clock observation or independent continuous-law cost sample. The complete individual service tables, including fitting warnings and repeated timings, appear in the technical supplement. Every earlier manuscript source and its labeled results remain available in exact retained copies and four compiled companions. The active reading path brings the constructive NBO, primitive theorem, query implementation and economic comparison together without deleting an unfavorable experiment or extending a theorem beyond its hypotheses.
